\documentclass[12pt, letterpaper]{hmcpset}
\usepackage{lin-alg-preamble}

\name{Jayden Thadani}
\class{Linear Algebra Done Right self-study}
\assignment{Exercises 6.B --- Orthonormal Bases}

\begin{document}

\begin{problem}[1]
	\begin{enumerate}
		\item
			Suppose $\theta \in \bb{R}$. Show that $(\cos{\theta}, \sin{\theta}),$ $(-\sin{\theta}, \cos{\theta})$ and $(\cos{\theta}, \sin{\theta}),$ $(\sin{\theta}, - \cos{\theta})$ are orthonormal bases of $\bb{R}^2$.
		\item
			Show that each orthonormal basis of $\bb{R}^2$ is of the form given by one of the two possibilities of part (a).
	\end{enumerate}
\end{problem}

\begin{solution}
	\begin{enumerate}
		\item
			We can show that both of the lists are orthonormal lists of the length $\dim \bb{R}^2 = 2$, and thus, bases.\\\\
			First, note that for $v = (\pm \cos{\theta}, \pm \sin{\theta})$
			\[
				v \cdot v = \cos^2{\theta} + \sin^2{\theta} = 1.
			\]
			Thus, for either of the lists $e_1, e_2$, we have that $e_j \cdot e_j = 1$ ($j \in \set{1, 2}$).\\\\
			Further, we can see that 
			\[
				(\cos{\theta}, \sin{\theta}) \cdot (- \sin{\theta}, \cos{\theta}) = \cos{\theta} \sin{\theta} - \cos{\theta} \sin{\theta} = 0
			\]
			and
			\[
				(\cos{\theta}, \sin{\theta}) \cdot (\sin{\theta}, -\cos{\theta}) = -(\cos{\theta}, \sin{\theta}) \cdot (- \sin{\theta}, \cos{\theta}) = 0.
			\]
			That is, for either of the lists, we have $e_1 \cdot e_2 = 0$.\\\\
			Thus, both of the lists are orthonormal. Since both of the lists are of length $2 = \dim \bb{R}^2$, they are also orthonormal bases.
		\item
			Suppose we have an orthonormal basis of $\bb{R}^2$, say, $e_1, e_2$.\\\\
			Then $e_1$ must be a unit vector in $\bb{R}^2$. For angle $\theta$ from the positive $x$-axis, we can write $e_1$ as $(\cos{\theta}, \sin{\theta})$. Since we are working with the Euclidean inner product, $e_2$ must then be $e_1$ rotated $\frac{\pi}{2}$ clockwise or counterclockwise:
			\[
				e_2 = (\cos{\theta + \frac{\pi}{2}}, \sin{\theta + \frac{\pi}{2}}) = (-\sin{\theta}, \cos{\theta})
			\]
			or
			\[
				e_2 = (\cos{\theta - \frac{\pi}{2}}, \sin{\theta - \frac{\pi}{2}}) = (\sin{\theta}, -\cos{\theta}).
			\]
			This gives us any arbitrary orthonomal basis of $\bb{R}^2$ in the desired form.
	\end{enumerate}
\end{solution}

\pagebreak

\begin{problem}[2]
	Suppose $e_1, \ldots, e_m$ is an orthonormal list of vectors in $V$. Let $v \in V$. Prove that
	\[
		\norm{v}^2 = \abs{\langle v, e_1 \rangle}^2 + \cdots + \abs{\langle v, e_m \rangle}^2
	\]
	if and only if $v \in \vspan{(e_1, \ldots, e_m)}$.
\end{problem}

\begin{solution}
	Suppose $v \in \vspan(e_1, \ldots, e_m)$. Then since $e_1, \ldots, e_m$ is an orthonormal basis of the vector space $\vspan{(e_1, \ldots, e_m)}$, we can write $v$ as
	\[
		v = \linn v, e_1 \rinn e_1 + \cdots + \linn v, e_m \rinn.
	\]
	This then naturally gives us
	\[
		\norm{v}^2 = \abs{\linn v, e_1 \rinn}^2 + \cdots + \abs{\linn v, e_m \rinn}^2.
	\]\\
	Suppose $v \notin \vspan{(e_1, \ldots, e_m)}$.\\\\
	Consider $u_{m + 1} = v - (\linn v, e_1 \rinn e_1 + \cdots + \linn v, e_m \rinn e_m) \neq 0$ and $e_{m + 1} = \frac{u_{m + 1}}{\norm{u_{m + 1}}}$. That is, expand the basis to include $v_m$ in its span. By the Gram-Schmidt procedure, $e_1, \ldots, e_{m + 1}$ is now an orthonormal basis such that $v \in \vspan (e_1, \ldots, e_{m + 1})$.\\\\
	This allows us to write 
	\[
		v = \linn v, e_1 \rinn e_1 + \cdots + \linn v, e_m \rinn e_m + \linn v, e_{m + 1} \rinn e_{m + 1}
	\]
	where $\linn v, e_{m + 1} \rinn \neq 0$ by our supposition. This gives us 
	\[
		\norm{v}^2 = \abs{\linn v, e_1 \rinn}^2 + \cdots + \abs{\linn v, e_m \rinn} + \abs{\linn v, e_{m + 1} \rinn}^2
	\]
	and from $\linn v, e_{m + 1} \rinn \neq 0$ it follows that
	\[
		\norm{v}^2 > \abs{\linn v, e_1 \rinn}^2 + \cdots + \abs{\linn v, e_m \rinn}^2
	\]
	and thus,
	\[
		\norm{v}^2 \neq \abs{\linn v, e_1 \rinn}^2 + \cdots + \abs{\linn v, e_m \rinn}^2.
	\]
\end{solution}

\pagebreak

\begin{problem}[3]
	Suppose $T \in \curly{L}(\bb{R}^3)$ has an upper-triangular matrix with respect to the basis $(1, 0, 0), (1, 1, 1), (1, 1, 2)$. Find an orthonormal basis of $\bb{R}^3$ (use the usual inner product on $\bb{R}^3$) with respect to which $T$ has an upper-triangular matrix.
\end{problem}

\begin{solution}
	\[
		\boxed{
			(1, 0, 0), (0, \frac{1}{\sqrt 2}, \frac{1}{\sqrt 2}), (0, -\frac{1}{\sqrt 2}, \frac{1}{\sqrt 2}).
		}
	\] \\
	Since the Gram-Schmidt process preserves upper-triangular matrices, we can just apply it to the given basis to get an orthonormal basis with respect to which $T$ has an upper-triangular matrix. Let the original basis be $v_1, v_2, v_3$ and the orthonormal basis be $e_1, e_2, e_3$. Let the non-normalised versions of the orthonormal basis be $u_1, u_2, u_3$.\\\\
	The Gram-Schmidt process gives us the following.
	\[
		e_1 = \frac{v_1}{\norm{v_1}} = \frac{(1, 0, 0)}{1} = (1, 0, 0).
	\]
	This allows us to compute $e_2$.
	\[
		\begin{split}
			u_2 & = v_2 - (v_2 \cdot e_1) e_1 \\
			    & = (1, 1, 1) - ((1, 1, 1) \cdot (1, 0, 0)) (1, 0, 0) \\
			    & = (1, 1, 1) - 1(1, 0, 0) \\ 
			    & = (0, 1, 1).
		\end{split}
	\]
	\[
		e_2 = \frac{u_2}{\norm{u_2}} = \frac{(0, 1, 1)}{\sqrt{2}} = (0, \frac{1}{\sqrt 2}, \frac{1}{\sqrt 2}).
	\]
	And then $e_3$.
	\[
		\begin{split}
			u_3 & = v_3 - (v_3 \cdot e_1)e_1 - (v_3 \cdot e_2)e_2 \\ 
			    & = (1, 1, 2) - ((1, 1, 2) \cdot (1, 0, 0))(1, 0, 0) - \frac{((1, 1, 2) \cdot (0, 1, 1))(0, 1, 1)}{\sqrt{2} \sqrt{2}} \\
			    & = (1, 1, 2) - 1(1, 0, 0) - \frac{3(0, 1, 1)}{2} \\
			    & = (0, -\frac{1}{2}, \frac{1}{2}).
		\end{split}
	\]
	\[
		e_3 = \frac{u_3}{\norm{u_3}} = \frac{(0, -\frac{1}{2}, \frac{1}{2})}{1 / \sqrt{2}} = (0, -\frac{1}{\sqrt 2}, \frac{1}{\sqrt 2}).
	\]\\
	Thus, $T$ has an upper-triangular matrix with respect to the orthonormal basis
	\[
		(1, 0, 0), (0, \frac{1}{\sqrt 2}, \frac{1}{\sqrt 2}), (0, -\frac{1}{\sqrt 2}, \frac{1}{\sqrt 2}).
	\]
\end{solution}

\pagebreak

\begin{problem}[4]
	Suppose $n$ is a positive integer. Prove that
	\[
		\frac{1}{\sqrt{2 \pi}}, \frac{\cos{x}}{\sqrt{\pi}}, \frac{\cos{2x}}{\sqrt{\pi}}, \cdots, \frac{\cos{nx}}{\sqrt{\pi}}, \frac{\sin{x}}{\sqrt{\pi}}, \frac{\sin{2x}}{\sqrt{\pi}}, \cdots, \frac{\sin{nx}}{\sqrt{\pi}}
	\]
	is an orthonormal list of vectors in $\curly{C}[-\pi, \pi]$, the vector space of continuous functions on $[-\pi, \pi]$ with inner product
	\[
		\linn f, g \rinn =  \int_{-\pi}^{\pi} f(x) g(x) \, dx.
	\]
\end{problem}

\begin{solution}
	We will show that with the inner product in the question
	\[
		\linn \frac{\cos{nx}}{\sqrt \pi}, \frac{\cos{nx}}{\sqrt \pi} \rinn = \linn \frac{\sin{nx}}{\sqrt \pi}, \frac{\sin{nx}}{\sqrt{\pi}} \rinn = \linn \frac{1}{\sqrt{2 \pi}}, \frac{1}{\sqrt{2 \pi}} \rinn = 1
	\]
	and
	\[
		\linn \frac{\cos{mx}}{\sqrt{\pi}} \frac{\cos{nx}}{\sqrt \pi} \rinn = \linn \frac{\sin{mx}}{\sqrt \pi}, \frac{\sin{nx}}{\sqrt \pi} \rinn = \linn \frac{1}{\sqrt{2 \pi}}, \sin{nx} \rinn = \linn \frac{1}{\sqrt{2 \pi}}, \cos{nx} \rinn = 0
	\]
	for $m \neq n$. This will give us that it is an orthonormal list.
	Finally,
	\[
		\linn \frac{\cos{mx}}{\sqrt{\pi}}, \frac{\sin{nx}}{\sqrt \pi} \rinn = 0
	\]
	regardless of whether $m = n$.\\\\
	For the first, we have the following calculations.
	\[
		\begin{split}
			\linn \frac{1}{\sqrt{2 \pi}}, \frac{1}{\sqrt{2 \pi}} \rinn & =\int_{-\pi}^\pi \Big ( \frac{1}{\sqrt{2 \pi}} \Big )^2 \, dx \\
										   & = \frac{\pi - (- \pi)}{2 \pi} \\ 
										   & = 1.
		\end{split}
	\]\\
	\[
		\begin{split}
			\linn \frac{\cos{nx}}{\sqrt \pi}, \frac{\cos{nx}}{\sqrt \pi} \rinn & = \frac{1}{\pi}\int_{-\pi}^\pi \cos^2{nx} \, dx \\
						       & = \frac{1}{2 \pi} \int_{-\pi}^\pi 1 + \cos{2nx} \, dx \\ 
						       & = \frac{1}{2 \pi} \Big ( x + \frac{\sin{2nx}}{2n} \Big ) \Big |_{-\pi}^{\pi} \\ 
						       & = \frac{1}{2 \pi} ((\pi + 0) - (-\pi + 0) ) \\
						       & = 1.
		\end{split}
	\]\\
	\[
		\begin{split}
			\linn \frac{\sin{nx}}{\sqrt \pi}, \frac{\sin{nx}}{\sqrt \pi} \rinn & = \frac{1}{\pi} \int_{-\pi}^{\pi} \sin^2{nx} \, dx \\
											   & = \frac{1}{2 \pi} \int_{- \pi}^\pi 1 - \cos{2nx} \, dx \\
											   & = \frac{1}{2 \pi} \Big ( x - \frac{\sin{2nx}}{2n} \Big ) \Big |_{-\pi}^\pi \\
											   & = \frac{1}{2 \pi} ((\pi - 0) - (-\pi - 0)) \\
											   & = 1.
		\end{split}
	\]\\
	For the second we have the following.
	\[
		\begin{split}
			\linn \frac{1}{\sqrt{2 \pi}}, \frac{\cos{nx}}{\sqrt{\pi}} \rinn  & = \frac{1}{\pi \sqrt 2}\int_{-\pi}^\pi \cos{nx} \, dx \\
											 & = \frac{1}{\pi \sqrt 2} \Big ( \frac{\sin{nx}}{n} \Big ) \Big |_{-\pi}^\pi \\
											 & = \frac{1}{\pi \sqrt 2} (0 - 0) \\
											 & = 0.
		\end{split}
	\]\\
	\[
		\begin{split}
			\linn \frac{1}{\sqrt{2 \pi}}, \frac{\sin{nx}}{\sqrt \pi} \rinn & = \frac{1}{\pi \sqrt 2} \int_{-\pi}^\pi \sin{nx} \, dx \\
										       & = \frac{1}{\pi \sqrt 2} \Big ( \frac{\cos{nx}}{n} \Big ) \Big |_{- \pi}^\pi \\
										       & = \begin{cases}
											       \frac{1}{\pi \sqrt 2} (\frac{1}{n} - \frac{1}{n}) & \text{for even } n \\
											       \frac{1}{\pi \sqrt 2} (\frac{-1}{n} - \frac{(-1)}{n}) & \text{for odd } n
										       \end{cases} \\
										       & = 0.
		\end{split}
	\]\\
	\[
		\linn \frac{\cos{mx}}{\sqrt \pi}, \frac{\cos{nx}}{\sqrt \pi} \rinn = \frac{1}{\pi} \int_{-\pi}^\pi \cos{mx} \cos{nx} \, dx.
	\]
	To compute the integral, we obtain an antiderivative by integrating by parts.
	\[
		\begin{split}
			\int \cos{mx} \cos{nx} \, dx & = \cos{mx} \int \cos{nx} \, dx - \int \Big ( \ddx{(\cos{mx})}{x} \int \cos{nx} \, dx \Big ) \, dx \\
						     & = \frac{\cos{mx} \sin{nx}}{n} - \int \frac{- m \sin{mx} \sin{nx}}{n} \, dx \\
						     & = \frac{\cos{mx \sin{nx}}}{n} + \frac{m}{n} \Big ( \sin{mx} \int \sin{nx} \, dx - \int \Big ( \ddx{(\sin{mx})}{x} \int \sin{nx} \, dx \Big ) \, dx \Big ) \\
						     & = \frac{\cos{mx} \sin{nx}}{n} + \frac{m}{n} \Big ( \frac{- \cos{nx} \sin{mx}}{n} - \int \frac{-m \cos{mx} \cos{nx}}{n} \, dx \Big ) \\
						     & = \frac{\cos{mx}\sin{nx}}{n} - \frac{m \cos{nx} \sin{mx}}{n^2} + \frac{m^2}{n^2} \int \cos{mx} \cos{nx} \, dx.
		\end{split}
	\]
	If we let $\curly I = \int \cos{mx} \cos{nx} \, dx$, then this gives us
	\[
		(1 - \frac{m^2}{n^2}) \curly I = \frac{n \cos{mx} \sin{nx} - m \cos{nx} \sin{mx}}{n^2}
	\]
	and since $m \neq n$, (and $m$ and $n$ are positive) $1 - \frac{m^2}{n^2} \neq 0$ and thus,
	\[
		\curly I = \frac{n \cos{mx} \sin{nx} - m \cos{nx} \sin{mx}}{n^2 - m^2}.
	\]
	Thus,
	\[
		\begin{split}
			\linn \frac{\cos{mx}}{\sqrt \pi}, \frac{\cos{nx}}{\sqrt \pi} \rinn & = \frac{1}{\pi} \Big ( \frac{n \cos{mx} \sin{nx} - m \cos{nx} \sin{mx}}{n^2 - m^2} \Big ) \Big |_{-\pi}^\pi \\
											   & = \frac{0 - 0}{\pi (n^2 - m^2)} \\
											   & = 0.
		\end{split}
	\]\\
	We could also have done the previous integral more simply by using a trigonometric identity, which we will now do for the case with $\sin$. Specifically, we will use the fact that $\sin{A} \sin{B} = \frac{\cos{(A - B)} - \cos{(A + B)}}{2}$
	\[
		\begin{split}
			\linn \frac{\sin{mx}}{\sqrt \pi}, \frac{\sin{nx}}{\sqrt \pi} \rinn & = \frac{1}{\pi} \int_{-\pi}^\pi \sin{mx} \sin{nx} \, dx \\
											   & = \frac{1}{2 \pi} \int_{-\pi}^\pi \cos{((m - n)x)} - \cos{((m + n)x)} \, dx \\
											   & = \frac{1}{2 \pi} \Big ( \frac{\sin{((m - n)x)}}{m - n} - \frac{\sin{((m + n)x)}}{m + n} \Big ) \Big |_{-\pi}^\pi \\
											   & = \frac{1}{2 \pi} (0 - 0) \\
											   & = 0.
		\end{split}
	\]\\
	Finally, for the third we have
	\[
		\begin{split}
			\linn \frac{\cos{mx}}{\sqrt \pi}, \frac{\sin{mx}}{\sqrt \pi} \rinn & = \frac{1}{\pi} \int_{-\pi}^\pi \cos{mx} \sin{nx} \, dx \\
											   & = \frac{1}{2 \pi} \int_{-\pi}^\pi \sin{((m + n)x)} - \sin{((m - n)x)} \, dx \\
											   & = \frac{1}{2 \pi} ( \cos{((m - n)x)} - \cos{((m + n)x)} ) \Big |_{-\pi}^\pi.
		\end{split}
	\]
	Since $m + n \equiv m - n \mod 2$, we have $\cos{((m + n)\pi)} = \cos{((m - n)\pi)}$. Then
	\[
		\begin{split}
			\frac{1}{2 \pi} ( \cos{((m - n)x)} - \cos{((m + n)x)} ) \Big |_{-\pi}^\pi & = \frac{1}{2 \pi}(\cos{((m + n) \pi)}) - \cos{((m + n) \pi)}) \\
												  & = 0.
		\end{split}
	\]
\end{solution}

\pagebreak

\begin{problem}[5]
	On $\curly{P}_2(\bb{R})$, consider the inner product given by
	\[
		\linn p, q \rinn = \int_0^1 p(x) q(x) \, dx.
	\]
	Apply the Gram-Schmidt procedure to the basis $1, x, x^2$ to produce an orthonormal basis of $\curly{P}_2(\bb{R})$
\end{problem}

\begin{solution}
	\[
		\boxed{
			1, 2 \sqrt 3 \Big ( x - \frac{1}{2} \Big ), 6 \sqrt 5 \Big ( x^2 - x + \frac{1}{6} \Big ).
		}
	\] \\
	Let $v_1, v_2, v_3$ denote the basis $1, x, x^2$, let $e_1, e_2, e_3$ denote the orthonormal basis obtained by applying the Gram-Schmidt procedure to it, and let $u_1, u_2, u_3$ denote the non-normalised version of $e_1, e_2, e_3$. Then
	\[
		e_1 = \frac{v_1}{\norm{v_1}} = \frac{1}{\int_0^1 \, dx} = 1
	\]
	This allows us to get $e_2$ as follows.
	\[
		\begin{split}
			u_2 & = v_2 - \linn v_2, e_1 \rinn e_1 \\
			    & = x - \int_0^1 x \, dx \\
			    & = x - \Big ( \frac{x^2}{2} \Big ) \Big |_0^1 \\
			    & = x - \frac{1}{2}.
		\end{split}
	\]
	\[
		\begin{split}
			\norm{u_2} & = \sqrt{\int_0^1 x^2 - x + \frac{1}{4} \, dx} \\
				   & = \sqrt{\Big ( \frac{x^3}{3} - \frac{x^2}{2} + \frac{x}{4} \Big ) \Big |_0^1} \\
				   & = \sqrt{\frac{1}{12}}
		\end{split}
	\]
	\[
		e_2 = \frac{u_2}{\norm{u_2}} \\ = \frac{x - \frac{1}{2}}{\sqrt{\frac{1}{12}}} = 2 \sqrt{3} \Big ( x - \frac{1}{2} \Big ).
	\]
	Finally, we can get $e_3$ as well.
	\[
		\begin{split}
			u_3 & = v_3 - \linn v_3, e_1 \rinn e_1 - \linn v_3, e_2 \rinn e_2 \\
			    & = x^2 - \int_0^1 x^2 \, dx - \sqrt{12} \, (x - \frac{1}{2}) \int_0^1 x^2 \sqrt{12} (x - \frac{1}{2}) \, dx \\
			    & = x^2 - \Big ( \frac{x^3}{3} \Big ) \Big |_0^1 - 12 \, (x - \frac{1}{2}) \Big ( \frac{x^4}{4} - \frac{x^3}{6} \Big ) \Big |_0^1 \\
			    & = x^2 - \frac{1}{3} - \Big ( x - \frac{1}{2} \Big ) \\
			    & = x^2 - x + \frac{1}{6}.
		\end{split}
	\]
	\[
		\begin{split}
			\norm{u_3} & = \sqrt{\int_0^1 x^4 - 2x^3 - \frac{4}{3} x^2 - \frac{1}{3} x + \frac{1}{36} \, dx} \\
				   & = \sqrt{\Big ( \frac{x^5}{5} - \frac{x^4}{2} - \frac{4 x^3}{9} - \frac{x^2}{6} + \frac{x}{36} \Big ) \Big |_0^1} \\
				   & = \sqrt{\frac{1}{180}}
		\end{split}
	\]
	\[
		e_3 = \frac{u_3}{\norm{u_3}} = 6 \sqrt{5} \Big (x^2 - x + \frac{1}{6} \Big )
	\]
\end{solution}

\pagebreak

\begin{problem}[6]
	Find an othonormal basis of $\curly{P}_2(\bb{R})$ (with inner product as in exercise 5) such that the differentiation operator (the operator that takes $p$ to $p'$) has an upper-triangular matrix with respect to this basis.
\end{problem}

\begin{solution}
	\[
		\boxed{
			1, 2 \sqrt 3 \Big ( x - \frac{1}{2} \Big ), 6 \sqrt 5 \Big ( x^2 - x + \frac{1}{6} \Big ).
		}
	\] \\
	The differentiation operator $D$ with $D p = p'$ has the upper-triangular matrix
	\[
		\begin{bmatrix}
			0 & 1 & 0 \\
			0 & 0 & 2 \\
			0 & 0 & 0
		\end{bmatrix}
	\]
	with respect to the basis $1, x, x^2$. \\\\
	Since the Gram-Schmidt preserves the spans of sub-lists of bases, it also preserves upper-triangular matrices. Thus, the basis obtained in problem 5, by applying the Gram-Schmidt process, gives us an upper triangular matrix with respect to it. Specifically, $D$ has the upper-triangular matrix
	\[
		\begin{bmatrix}
			0 & 2 \sqrt 3 & 0 \\
			0 & 0 & 2 \sqrt 3 \sqrt 5 \\
			0 & 0 & 0
		\end{bmatrix}
	\]
	with respect to the orthonormal basis $1, 2 \sqrt 3 \Big ( x - \frac{1}{2} \Big ), 6 \sqrt 5 \Big ( x^2 - x + \frac{1}{6} \Big )$.
\end{solution}

\pagebreak

\begin{problem}[7]
	Find a polynomial $q \in \curly{P}_2(\bb{R})$ such that
	\[
		p(\frac{1}{2}) = \int_0^1 p(x) q(x) \, dx
	\]
	for every $p \in \curly{P}_2(\bb{R})$.
\end{problem}

\begin{solution}
	\[
		\boxed{
			q = -15 (x^2 - x) - \frac{3}{2}.
		}
	\] \\
	Consider the inner-product on $\mathcal{P}(\bb{R}^2)$ from problems 5 and 6 and the linear functional $\varphi \in \mathcal{P}(\bb{R}^2)^\vee$ given by $\varphi(p) = p(\frac{1}{2})$. Then we can rephrase this problem as finding some $q$ such that
	\[
		\linn p, q \rinn = \varphi(p).
	\]\\
	Then we know from the Riesz representation theorem that, given an orthonormal basis $e_1, e_2, e_3$, (which we computed in problem 5) we have
	\[
		\begin{split}
			q & = \overline{\varphi(e_1)} \, e_1 + \overline{\varphi(e_2)} \, e_2 + \overline{\varphi(e_3)} \, e_3 \\
			  & = e_1 + 2 \sqrt 3 \, \Big ( \frac{1}{2} - \frac{1}{2} \Big ) e_2 + 6 \sqrt 5 \, \Big ( \frac{1}{4} - \frac{1}{2} + \frac{1} {6}  \Big ) e_3 \\
			  & = 1 - \frac{6 \sqrt 5}{12} \, 6 \sqrt 5 \, \Big ( x^2 - x + \frac{1}{6} \Big ) \\
			  & = -15 (x^2 - x) - \frac{3}{2}.
		\end{split}
	\]
	
\end{solution}

\pagebreak

\begin{problem}[8]
	Find a polynomial $q$ such that
	\[
		\int_0^1 p(x) \cos{\pi x} \, dx = \int_0^1 p(x) q(x) \, dx
	\]
	for every $p \in \curly{P}_2(\bb{R})$.
\end{problem}

\begin{solution}
	\[
		\boxed{
			q = - \frac{24}{\pi^2} \left ( x - \frac{1}{2} \right ).
		}
	\] \\
	Consider the functional on $\varphi \in \mathcal P_2(\mathbb R)^\vee$ by
	\[
		\varphi : p \mapsto \int_0^1 p(x) \cos{\pi x} \, dx.
	\]
	We know from the Riesz representation theorem that
	\[
		q = \overline{\varphi(e_1)} \, e_1 + \overline{\varphi(e_2)} \, e_2 + \overline{\varphi(e_3)} \, e_3
	\]
	for the orthonormal basis that we calculated in problem 5. \\\\
	Thus, we can compute
	\[
		\begin{split}
			q & = \overline{ \left ( \int_0^1 e_1 \, \cos{\pi x} \, dx \right )} \, e_1 + \overline{\left ( \int_0^1 e_2 \, \cos{\pi x} \, dx \right )} \, e_2 + \overline{\left ( \int_0^1 e_3 \, \cos{\pi x} \, dx \right )} \, e_3 \\
			  & = 0 \, e_1 - \frac{4 \sqrt 3}{\pi^2} \, e_2 + 0 \, e_3 \\
			  & = - \frac{24}{\pi^2} \left ( x - \frac{1}{2} \right )
		\end{split}
	\]
\end{solution}

\pagebreak

\begin{problem}[9]
	What happens if the Gram-Schmidt procedure is applied to a list of vectors that is not linearly independent?
\end{problem}

\begin{solution}
	Suppose $v_1, \dots, v_m$ is linearly dependent. Thus, there must be some least $j \in \mathbb N_m$ such that $v_j \in \operatorname{span}(v_1, \dots, v_{j - 1})$. Note then that the Gram-Schmidt procedure works normally for $v_1, \dots, v_{j - 1}$, producing the orthonormal list $e_1, \dots, e_{j - 1}$ with the same span. Thus, $v_j \in \operatorname{span}(e_1, \dots, e_{j - 1})$. From here, we notice that $v_j = \left < v_j, e_1 \right > e_1 + \dots + \left < v, e_{j - 1} \right > e_{j - 1}$, and thus, that $e_j$, which is just a scaled version of $v - v_j = \left < v_j, e_1 \right > e_1 + \dots + \left < v, e_{j - 1} \right > e_{j - 1}$ must be $0$. That is, if we apply the Gram-Schmidt a linearly dependent list, the first vector in the span of the vectors before it will give $0$.
\end{solution}

\pagebreak

\begin{problem}[10]
	Suppose $V$ is a real inner product space and $v_1 \ldots v_m$ is a linearly independent list of vectors in $V$. Prove that there exist exactly $2^m$ orthonormal lists $e_1, \ldots, e_m$ of vectors in $V$ such that
	\[
		\vspan (v_1, \ldots, v_j) = \vspan (e_1, \ldots, e_j)
	\]
	for all $j \in \bb{N}_m$
\end{problem}

\begin{solution}
	Apply the Gram-Schmidt procedure to $v_1, \dots, v_m$ to get an orthonormal basis $e_1, \dots, e_m$. Choose to multiply any subset of them by $-1$. This preserves the span of each sublist $e_1, \dots, e_j$. There are $2^m$ ways to do this. \\\\
	We also have that $e_{j + 1}$ and $-e_{j + 1}$ are the only way to add $v_{j + 1}$ to the span of $e_1, \dots, e_j$ while preserving the orthonormal property. We can prove this as follows. \\\\
	Suppose we have $\operatorname{span}(e_1, \dots e_j) = \operatorname{span}(v_1, \dots, v_j)$ and we want $e_{j + 1}$ such that $e_1, \dots, e_{j + 1}$ is orthonormal, and $\operatorname{span}(e_1, \dots, e_{j + 1}) = \operatorname{span}(v_1, \dots, v_{j + 1}) = \operatorname{span}(e_1, \dots, e_j, v_{j + 1})$. Then notice that we can write $e_{j + 1} \in \operatorname{span}(e_1, \dots, e_j, v_{j + 1})$ as
	\[
		e_{j + 1} = c_1 e_1 + \dots + c_j e_j + d v_{j + 1}.
	\]
	Taking the inner product of everything with $e_{k}$ (for some $k < j$) gives us
	\[
		0 = 0 + \dots + c_k + \dots + 0 + d \left < v_{j + 1}, e_k \right >.
	\]
	Thus, all $c_k = - d \left < v_{j + 1}, e_k \right >$. This allows us to write $e_{j + 1}$ as
	\[
		e_{j + 1} = d ( v_{j + 1} - ( \left < v_{j + 1}, e_1 \right > e_1 + \dots + \left < v_{j + 1}, e_j \right > e_j ) )
	\]
	Finally, using our fact that $\Vert e_{j + 1} \Vert^2 = 1$, we take the norm on both sides and get
	\[
		1 = d^2 \left ( \Vert v_{j + 1} \Vert^2 - \sum_{i = 1}^j \left < v_{j + 1}, e_i \right >^2 \right ).
	\]
	This gives us two options for $d$:
	\[
		d = \pm \frac{1}{\sqrt{\left ( \Vert v_{j + 1} \Vert^2 - \sum_{i = 1}^j \left < v_{j + 1}, e_i \right >^2 \right )}}.
	\]
	Since $d$ uniquely determines all of the coefficients, and the $e_{j + 1}$ and $- e_{j + 1}$ we get from the Gram-Schmidt process are valid choices, they must be the only two.
\end{solution}

\pagebreak

\begin{problem}[11.]
	Suppose $\left < \cdot, \cdot \right >_1$ and $\left < \cdot, \cdot \right >_2$ are inner products on $V$ such that $\left < v, w \right >_1 = 0$ if and only if $\left < v, w \right >_2 = 0$. Prove that there is a positive number $c$ such that $\left < v, w \right >_1 = c \left < v, w \right >_2$ for every $v, w \in V$.
\end{problem}

\begin{solution}
	Suppose $v \in V$ gives us
	\[
		\frac{\Vert v \Vert_1^2}{\Vert v \Vert_2^2} = c_v > 0.
	\] \\
	Notice that we can write any $w \in V$ in terms of $v$. That is, we can get the orthogonal decompositions
	\[
		w = a_i v + w'
	\]
	where
	\[
		a_i = \frac{\left < w, v \right >_i}{\Vert v \Vert_i^2}
	\] \\
	Note that since we have $\langle v', u \rangle_i = 0$ we have
	\[
		\begin{split}
			\left < v, w - a_1 v \right >_1 & = 0 \\
			\left < v, w - a_1 v \right >_2 & = 0 \\
			\left < v, w - \frac{\langle w, v \rangle_1}{\Vert v \Vert_1^2} v \right >_2 & = 0 \\
			\langle v, w \rangle_2 - \overline{\langle w, v \rangle_1} \frac{\Vert v \Vert_2^2}{\Vert u \Vert_1^2} & = 0 \\
			\langle v, w \rangle_2 - \frac{1}{c_v} \langle v, w \rangle_1 & = 0
		\end{split}
	\]
	and thus
	\[
		\langle v, w \rangle_1 = c_v \langle v, w \rangle_2.
	\]
	Now we just need to show that $c_u = c_v$ for any $u \in V$. Note that for any $u \in V$ we have
	\[
		\langle v, u \rangle_1 = c_v \langle v, u \rangle_2.
	\]
	but also
	\[
		\langle u, v \rangle_1 = c_u \langle u, v \rangle_2.
	\]
	Note that since $c_u \in \mathbb{R}$, we can rewrite this as 
	\[
		\langle v, u \rangle_1 = c_u \langle v, u \rangle_2
	\]
	by taking conjugates on both sides. Then
	\[
		(c_v - c_u) \langle v, u \rangle_2 = 0
	\]
	for all $u \in U$. This means that whenever, $u$ and $v$ are not orthogonal, we have $c_u = c_v$. When they are orthogonal, consider $v + u$. We have $c_v = c_{v + u}$ and $c_u = c_{v + u}$ since neither of them is orthogonal to $v + u$, and thus, $c_v = c_u$. Let $c = c_v$ for all $v \in V$. \\\\
	Thus, for all $v, w \in V$ we have $\langle v, w \rangle_1 = c \langle v, w \rangle_2$.
\end{solution}

\pagebreak

\begin{problem}[12.]
	Suppose $V$ is finite-dimensional and $\left < \cdot, \cdot \right >_1, \left < \cdot, \cdot \right >_2$ are inner products on $V$ with corresponding norms $\lVert \cdot \rVert_1$ and $\lVert \cdot \rVert_2$. Prove that there exists a positive number $c$ such that
	\[
		\lVert v \rVert_1 \le c \lVert v \rVert_2
	\]
	for every $v \in V$.
\end{problem}

\begin{solution}
	Let $v_1, \dots, v_n$ be a unit basis of $V$. Let $e_1, \dots, e_n$ be the orthonormal basis obtained by applying the Gram-Schmidt procedure with respect to the first inner product. Let  $f_1, \dots, f_n$ be the orthonormal basis obtained by applying the Gram-Schmidt procedure with respect to the second inner product. \\\\
	For any $v = a_{1} v_{1} + \dots + a_n v_{n}$, we can also write $v = b_{1} e_{1} + \dots + b_{n} e_{n}$ and $v = c_{1} f_{1} + \dots + c_{n} f_{n}$. We claim that we have $\lVert v \rVert_{1}^{2} \le \sum_{i = 1}^{n} a_{i}^{2}$. \\\\
\end{solution}

\pagebreak

\begin{problem}[13.]
	Suppose $v_1, \dots, v_m$ is a linearly independent list in $V$. Show that there exists $w \in W$ such that $\left < w, v_j \right > > 0$ for all $j \in \mathbb N_m$.
\end{problem}

\pagebreak

\begin{problem}[14.]
	Suppose $e_1, \dots, e_n$ is an orthonormal basis of $V$ and $v_1, \dots, v_n$ are vectors in $V$ such that
	\[
		\Vert e_j - v_j \Vert < \frac{1}{\sqrt n}
	\]
	for each $j \in \mathbb N_n$. Prove that $v_1, \dots, v_n$ is a basis of $V$.
\end{problem}

\pagebreak

\begin{problem}[15.]
	Suppose $\mathcal C_{\mathbb R}([-1, 1])$ is the vector space of continuous real-valued functions on the interval $[-1, 1]$ with inner product given by
	\[
		\left < f, g \right > = \int_{-1}^1 f(x) g(x) \, dx
	\]
	for $f, g \in \mathcal C_{\mathbb R}([-1, 1])$. Let $\varphi$ be the linear functional on $\mathcal C_{\mathbb R}([-1, 1])$ defined by $\varphi(f) = f(0)$. Show that there does not exist $g \in \mathcal C_{\mathbb R}([-1, 1])$ such that
	\[
		\varphi(f) = \left < f, g \right >
	\]
	for every $f \in \mathcal C_{\mathbb R}([-1, 1])$.
\end{problem}

\pagebreak

\begin{problem}[16.]
	Suppose $V$ is a finite-dimensional vector space over $\mathbb C$, $T \in \mathcal L(V)$, all the eigenvalues of $T$ have absolute value less than $1$, and $\varepsilon > 0$. Prove that there exists a positive integer $m$ such that $\Vert T^m v \Vert \le \varepsilon \Vert v \Vert$ for every $v \in V$.
\end{problem}

\begin{solution}~
	
	Since  $V$ is a complex vector space, $T$ must have an upper-triangular matrix $A$ with respect to an orthonormal basis $e_{1}, \dots, e_{n}$. Note that since all the eigenvalues are less than $1$, we have $\lvert A_{j, j} \rvert < 1$ for all $j \in \mathbb{N}_{n}$. \\

	Then, for any $v = a_{1} e_{1} + \cdots + a_{n} e_{n}$, we have
	\[
		Tv = \sum_{j = 1}^{n} \left ( \sum_{k = j}^{n} a_{k} A_{j, k} \right ) e_{j}.
	\]
	We want to bound the magnitudes of the coefficients $\lvert \sum_{k = j}^{n} a_{k} A_{j, k} \rvert$ to be less than or equal to $\lambda_{j} \lvert a_{j} \rvert$ where $0 < \lambda_j < 1$. Unfortunately, with $A$, this is only possible for $j = n$ where the coefficient is $\lvert A_{n, n} a_{n} \rvert \le \lvert A_{n, n} \rvert \lvert a_{n} \rvert$ --- in all other cases, the rest of the terms $a_{k} A_{j, k}$ can be too large. \\

	However, by scaling our basis, we can make these terms arbitrarily small. Choose some positive $\delta < 1$, and then set $f_{n + 1 - j} = \delta^{j} e_{n + 1 - j}$ for all $j \in \mathbb{N}_{n}$. Let $B$ be the upper-triangular matrix of $T$ with respect to $f_{1}, \dots, f_{n}$. Then $B_{i, j} = \delta^{j - i} A_{i, j}$. Moreover, for any $v = b_{1} f_{1} + \cdots + b_{n} f_{n} = a_{1} e_{1} + \cdots + a_{n} e_{n}$ we have $b_{j} = a_{j}/\delta^{n + 1 - j}$ \\

	This means that, for any such $v$, we have
	\[
		Tv = \sum_{j = 1}^{n} \left ( \sum_{k = j}^{n} b_{k} B_{j, k} \right ) e_{j}
	\]
	just as before, but also, we have the bound
	\[
		\begin{split}
			\left \lvert \sum_{k = j}^{n} b_{k} B_{j, k} \right \rvert & \le \left \lvert b_{j} B_{j, j} \right \rvert + \sum_{k = j + 1}^{n} \left \lvert b_{k} B_{j, k} \right \rvert \\
										   & \le \lvert b_{j} \rvert \lvert A_{j, j} \rvert + \delta \sum_{k = j + 1}^{n} \lvert a_{j} A_{j, k} \rvert \\
										   & \le \lambda_{j} \lvert b_{j} \rvert
		\end{split}
	\]
	for any $\lambda_{j}$ with $\lvert A_{j, j} \rvert < \lambda_{j} < 1$. \\

	Finally, we can choose $\lambda = \max_{j \in \mathbb{N}_{n}} \lambda_{j}$ and get  $\lVert Tv \rVert \le \lambda \lVert v \rVert$. Thus, for any positive integer $m \ge \log_{\lambda} \varepsilon$, we have $\lVert T^{m} v \rVert \le \varepsilon \lVert v \rVert$ for all $v \in V$.
\end{solution}

\pagebreak

\begin{problem}[17.]
	For $u \in V$, let $\Phi u$ denote the functional on $V$ defined by
	\[
		(\Phi u)(v) = \left < v, u \right >
	\]
	for $v \in V$.
	\begin{enumerate}
		\item
			Show that if $\mathbb F = \mathbb R$, then $\Phi$ is a linear map from $V$ to $V^{\vee}$.
		\item
			Show that if $\mathbb F = \mathbb C$, and $V \neq \{ 0 \}$, then $\Phi$ is not a linear map.
		\item
			Show that $\Phi$ is injective.
		\item
			Suppose $\mathbb F = \mathbb R$ and $V$ is finite-dimensional. Use parts $(a)$ and $(c)$ and a dimension-counting argument to show that $\Phi$ is an isomorphism from $V$ onto $V^{\vee}$.
	\end{enumerate}
\end{problem}

\begin{solution}~
	\begin{enumerate}
		\item Suppose $u = \lambda u_{1} + u_{2}$ for some $\lambda \in \mathbb{R}$. Then, for any $v \in V$, we would have
			\[
				\begin{split}
					(\Phi u)(v) & = \left < v, u \right > \\
						    & = \left < v, \lambda u_{1} + u_{2} \right > \\
						    & = \left < v, \lambda u_{1} \right > + \left < v, u_{2} \right > \\
						    & = \lambda \left < v, u_{1} \right > + \left < v, u_{2} \right > \\
						    & = \lambda (\Phi u_{1})(v) + (\Phi u_{2})(v) \\
						    & = (\lambda \Phi u_{1} + \Phi u_{2})(v).
				\end{split}
			\]
			Thus, $\Phi (\lambda u_{1} + u_{2}) = \lambda \Phi u_{1} + \Phi u_{2}$ --- $\Phi$ is linear.
		\item Suppose $V \neq \{ 0 \}$. Then we have some $u, v \in V$ with $u \neq 0$. Notice that
			\[
				\begin{split}
					(\Phi (iu))(v) & = \left < v, iu \right > \\
						    & = -i \left < v, u \right > \\
						    & = - i (\Phi u)(v)
				\end{split}
			\]
			We can only have linearity if $(\Phi u)(v) = 0$ always. However, if we choose $v = u \neq 0$, (which we can do) then $\Phi(u) \neq 0$. Thus, $\Phi$ is not linear.
		\item Suppose we have $u_{1}, u_{2} \in V$ with $\Phi u_{1} = \Phi u_{2}$. Then we must have $(\Phi u_{1})(v) = (\Phi u_{2})(v)$ at each $v \in V$, and thus, $\left < v, u_{1} \right > = \left < v, u_{2} \right >$ at each $v \in V$. But then, by additivity in the last slot, we have $\left < v, u_{1} - u_{2} \right > = 0$ always, even for $v = u_{1} - u_{2}$. Thus, $\left < u_{1} - u_{2}, u_{1} - u_{2} \right > = 0$. This gives us $u_{1} - u_{2} = 0$, and finally, $u_{1} = u_{2}$.
		\item We already know that $\dim V = \dim V^{\vee}$. Thus, $\Phi$, an injective linear map from $V$ to $V^{\vee}$ is an isomorphism.
	\end{enumerate}
\end{solution}

\end{document}
